% Part: methods
% Chapter: induction
% Section: structural-induction

\documentclass[../../../include/open-logic-section]{subfiles}

\begin{document}

\olfileid{mth}{ind}{sti}

\olsection{جوړښتي استقرا}

تر اوسه مو استقرا د ټولو طبيعي عددونو په هکله د
نتيجو د ثابتولو دپاره کارولې ده۔ خو يو ورته اصل
مستقيم د يو په استقرايي ډول تعريف شوي سټ د ټولو
!!{element}s په هکله د نتيجو د ثابتولو دپاره
کارېداے شي۔ دې ته ډېر ځله \emph{جوړښتي} استقرا
وايي، ځکه چې د په استقرايي ډول تعريف شوو څيزونو
پر جوړښت تکيه کوي۔

په عمومي ډول، يو استقرايي تعريف په (الف)~د سټ د
«لومړنيو» !!{element}s په لړليک او (ب)~د داسې عملونو
په لړليک ورکول کېږي چې د سټ له پخوانيو غړو
څخه ئې نوي !!{element}s جوړوي۔ د بېلګې په توګه، د
ښه‌جوړو ترمونو په حالت کښې لومړني څيزونه تورې دي۔
يوازې يو عمل لرو؛ عملونه دا دي:
\begin{align*}
  o(s_1, s_2) = & [s_1 \circ s_2]
\end{align*}
آن طبيعي عددونه~$\Nat$ پخپله هم د يو استقرايي
تعريف په ذريعه ورکړل شوي ګڼلاے شئ: لومړنے څيز~$0$
دے، او عمل د راتلونکي تابع~$x + 1$ ده۔

د يو په استقرايي ډول تعريف شوي سټ د ټولو غړو په
هکله د څه ثابتولو دپاره، يعنې دا چې د سټ هر
!!{element} خاصيت~$P$ لري، بايد:
\begin{enumerate}
\item ثابت کړو چې لومړني څيزونه~$P$ لري۔
\item ثابت کړو چې د هر عمل~$o$ دپاره، که ارګومېنټونه
  ئې~$P$ ولري، نو نتيجه ئې هم لري۔
\end{enumerate}
د بېلګې په توګه، د ټولو ښه‌جوړو ترمونو په هکله
د څه ثابتولو دپاره به وښيو چې د ټولو تورو دپاره
رښتيا دي، او د $[s_1 \circ s_2]$ دپاره هم رښتيا
دي، په دې شرط چې د $s_1$ او $s_2$ هر يوه دپاره
په جلا توګه رښتيا وي۔

\begin{prop}
د $[$ شمېر د $]$ له شمېر سره برابر دے، په هر
ښه‌جوړ ترم~$t$ کښې۔
\end{prop}

\begin{proof}
جوړښتي استقرا کاروو۔ ښه‌جوړ ترمونه په استقرايي
ډول تعريف شوي، چې تورې ئې لومړني څيزونه دي او
عمل $o$ ئې له پخوانيو څخه نوي ښه‌جوړ ترمونه جوړوي۔
\begin{enumerate}
\item دعوه د هر توري دپاره رښتيا ده، ځکه چې د $[$
  شمېر په يوازې يوه توري کښې~$0$ دے او د $]$ شمېر
  پکښې هم~$0$ دے۔
\item فرض کړئ چې د $[$ شمېر په $s_1$ کښې د $]$
  له شمېر سره برابر دے، او همدا خبره د $s_2$
  دپاره هم رښتيا ده۔ د $[$ شمېر په $o(s_1, s_2)$
  کښې، يعنې په $[s_1 \circ s_2]$ کښې، د $[$ د
  شمېرونو مجموعه ده په $s_1$ او $s_2$ کښې، له
  يو سره۔ د $]$ شمېر په $o(s_1, s_2)$ کښې د $]$
  د شمېرونو مجموعه ده په $s_1$ او $s_2$ کښې، له
  يو سره۔ نو د $[$ شمېر په $o(s_1, s_2)$ کښې د
  $]$ له شمېر سره برابر دے په $o(s_1,s_2)$ کښې۔
\end{enumerate}
\end{proof}

\begin{prob}
په جوړښتي استقرا ثابت کړئ چې هېڅ ښه‌جوړ ترم
په~$]$ نۀ پيلېږي۔
\end{prob}

په جوړښتي استقرا يو بل ثبوت ورکوو: د نښو د
يو تار~$t$ خاصه پيلنۍ برخه هر هغه تار~$s$ دے
چې له کيڼې خوا په لوستلو د نښې په نښه له $t$
سره يو شان وي، خو $t$ اوږد وي۔ نو، د بېلګې په
توګه، $[a \circ {}$ خاصه پيلنۍ برخه ده د
$[a \circ b]$؛ خو نۀ $[b \circ {}$، ځکه چې په
دويمه نښه کښې توپير لري، او نۀ $[a
\circ b]$، ځکه چې دواړه يو شان اوږد دي۔

\begin{prop}\ollabel{prop:initial}
د يوه ښه‌جوړ ترم~$t$ هره خاصه پيلنۍ برخه
د $[$ نښې د $]$ له نښو زياتې لري۔
\end{prop}
\textbf{د سرچينې يادونه OLSTH-042:} د دې دعوې
او لاندې ثبوت دپاره د ناتشې پيلنۍ برخې شرط
پکار دے۔ پورته چاپي تعريف تش تار نۀ باسي؛
که تش تار هم د خاصو پيلنيو برخو په شمېر کښې
راشي، د دواړو قوسونو شمېر ئې صفر دے او
دا سخت زياتوالے نۀ پوره کوي۔ د يوه توري
تش پيل هم همدا استثنا ده۔ چاپي دعوه او
ثبوت ساتل شوي؛ د شرط دغه نشتوالے دلته
په ښکاره ياد شو۔

\begin{proof}
پر~$t$ په استقرا:
  \begin{enumerate}
  \item $t$ يوازې يو تورے دے: بيا $t$ هېڅ خاصه پيلنۍ برخه نۀ لري۔
  \item $t = [s_1 \circ s_2]$، د ځينو ښه‌جوړو ترمونو
  $s_1$ او~$s_2$ دپاره۔ که $r$ د $t$ يوه خاصه
  پيلنۍ برخه وي، نو څو امکانات شته:
    \begin{enumerate}
    \item $r$ يوازې $[$ دے: بيا $r$ يو $[$ زيات لري
      له خپلو~$]$ څخه۔
    \item $r$، $[r_1$ دے، چې $r_1$ خاصه پيلنۍ برخه
      ده د~$s_1$۔ ځکه چې $s_1$ ښه‌جوړ ترم دے، د
      استقرايي فرض له مخې $r_1$ د $[$ نښې له
      $]$ څخه زياتې لري، او همدا خبره د $[r_1$
      دپاره هم رښتيا ده۔
    \item $r$، $[s_1$ يا $[s_1 \circ {}$ دے: د
      تېرې نتيجې له مخې د $[$ او $]$ شمېرونه
      په~$s_1$ کښې برابر دي؛ نو د $[$ شمېر
      په $[s_1$ يا $[s_1 \circ {}$ کښې يو زيات
      دے د~$]$ له شمېر څخه۔
    \item $r$، $[s_1 \circ r_2$ دے، چې $r_2$ خاصه
      پيلنۍ برخه ده د~$s_2$۔ د استقرايي فرض
      له مخې $r_2$ د $[$ نښې له $]$ څخه زياتې
      لري۔ د تېرې نتيجې له مخې د $[$ او
      $]$ شمېرونه په~$s_1$ کښې برابر دي۔ نو
      د $[$ شمېر په $[s_1 \circ r_2$ کښې زيات
      دے د~$]$ له شمېر څخه۔
    \item $r$، $[s_1 \circ s_2$ دے: د تېرې نتيجې
      له مخې د $[$ او $]$ شمېرونه په $s_1$
      کښې برابر دي، او همدا د~$s_2$ دپاره
      هم ده۔ نو يو $[$ زيات دے په $[s_1 \circ s_2$
      کښې، د~$]$ له شمېر څخه۔
    \end{enumerate}
  \end{enumerate}
\end{proof}

\end{document}
